\section[Stage 056]{Stage 056: Transport Origin of Lowest-Lane Source Asymmetry}
\label{stage:056}

\stagefield{Part / anchor}{Part III; secondary anchor \resultanchor{MTDC-T7}.}
\claimstatus{\StatusExactClosure{} for the drift-diffusion zero-flux source family.}
\stagefield{Purpose}{Stage~056 gives the exponential source profile a physical transport origin rather than treating it as a free shape ansatz.}
\stagefield{Inputs}{A one-dimensional source density \(\sigma(s)\) on the finite throat and a drift-diffusion current.}

\stagefield{Verification}{SymPy audit: \StageFile{scripts/moving_throat_pde_stage056_transport_source_asymmetry_sympy_audit.py}; transcript: \StageFile{scripts/output/moving_throat_pde_stage056_transport_source_asymmetry_sympy_audit.txt}.  Mathematica audit: \StageFile{mathematica/moving_throat_pde_stage056_transport_source_asymmetry_mathematica_audit.wl}; transcript: \StageFile{mathematica/output/moving_throat_pde_stage056_transport_source_asymmetry_mathematica_audit.txt}.}

\subsection*{Derivation}

Take
\begin{equation}
\label{eq:app-stage056-transport}
\partial_t\sigma+\partial_sJ=0,
\qquad
J=-D_\sigma\partial_s\sigma+v_\sigma\sigma.
\end{equation}
The stationary zero-flux equation \(J=0\) gives
\begin{equation}
\label{eq:app-stage056-exp-profile}
\sigma(s)\propto e^{v_\sigma s/D_\sigma}.
\end{equation}
With
\begin{equation}
\label{eq:app-stage056-Pe}
\boxed{
\mathrm{Pe}=\frac{v_\sigma L}{D_\sigma}},
\end{equation}
the normalized stationary family is exactly the Stage~053 exponential source with \(\alpha=\mathrm{Pe}\).  Therefore
\begin{equation}
\label{eq:app-stage056-Omega-Pe}
\boxed{
\Omega_{\mathrm{Pe}}
=\frac{\pi\mathrm{Pe}(2\mathrm{Pe}e^{\mathrm{Pe}}+\pi)}{(4\mathrm{Pe}^2+\pi^2)(e^{\mathrm{Pe}}-1)}}.
\end{equation}

\stagefield{Output}{Transport interpretation \eqref{eq:app-stage056-Pe} and physical overlap factor \eqref{eq:app-stage056-Omega-Pe}.}
\stagefield{Downstream use}{Stage~057 uses \(\mathrm{Pe}\) as the physical transport coordinate of \(\zeta_0\).}

